\documentclass{report}
\usepackage{amsmath,amssymb}
\setlength{\parindent}{0mm}
\setlength{\parskip}{1em}
\begin{document}
\begin{center}
\rule{6in}{1pt} \
{\large Tzanio Kolev \\
{\bf Parallel Auxiliary Space AMG Solver for H(Div) Problems}}

Center for Applied Scientific Computing \\ Lawrence Livermore National Laboratory \\ P O Box 808 \\ L-560 \\ Livermore \\ CA 94551
\\
{\tt tzanio@llnl.gov}\\
Panayot Vassilevski\end{center}

In this talk we present a family of scalable preconditioners for
matrices arising in the discretization of H(div) problems using
the lowest order Raviart-Thomas finite elements. Our approach
belongs to the class of ``auxiliary space''-based methods and
requires only the finite element stiffness matrix plus some
minimal additional discretization information about the topology
and orientation of the mesh entities. We outline the theory
which is based on certain regular decompositions for H(div)
spaces by relating them to a previously studied ``inf-sup''
condition for parameter-dependent Stokes problems. We provide a
detailed algebraic description of the method, its parallel
implementation and different variants of this parallel auxiliary
space divergence solver. Also, we discuss its relations to the
Hiptmair-Xu auxiliary space decomposition of Raviart-Thomas
spaces [1], the related approaches [4,5], as well as the
auxiliary space Maxwell solver (AMS) [2] from the hypre library
[3]. An extensive set of numerical experiments demonstrate the
robustness and scalability of our implementation on large-scale
parallel H(div) problems with large jumps in the material
coefficients.

[1] R.Hiptmair and J.Xu, Nodal auxiliary space preconditioning in H(curl)
and H(div) spaces, SIAM J. Num. Anal., 45 (2007), pp. 2483-2509.

[2] Tz. Kolev and P. Vassilevski, Parallel Auxiliary Space AMG for
H(curl) Problems, special issue on ``Adaptive and Multilevel Methods in
Electromagnetics'', Journal of Computational Mathematics, (27) 2009, pp.
604-623.

[3] hypre}: a library of high performance preconditioners,
http://www.llnl.gov/CASC/hypre/

[4] R.Tuminaro, J.Xu, and Y.Zhu, Auxiliary space preconditioners for
mixed finite element methods, in Domain Decomposition Methods in Science
and Engineering XVIII, vol.70 of Lecture Notes in Computational Science
and Engineering, Springer, 2009, pp. 99-109.

[5] P.Bochev, C.Siefert, R.Tuminaro, J.Xu, and Y.Zhu, Compatible gauge
approaches for H(div) equations, Tech. Rep. SAND 2007-5384P, Sandia
National Laboratories, 2007.


\end{document}
